\documentclass[10pt,a4paper]{amsart}
\usepackage[margin=19mm]{geometry}
\usepackage{amsmath,amssymb,mathtools}
\usepackage[scr=rsfs,cal=euler]{mathalfa}
\usepackage{xfrac}
\usepackage[unicode,hidelinks,bookmarks=false]{hyperref}
\newcommand{\ie}{i.e.\ }
\newcommand{\cf}{cf.\ }
\newcommand{\resp}{respectively\ }
\newcommand{\Subsection}{\subsection}
\providecommand{\nobreakdash}{\nobreak\hbox{-}}
\setlength{\emergencystretch}{2em}
\multlinegap0pt
\usepackage{amssymb}
\usepackage{mathtools}
\newtheorem{lemma}{Lemma}
\DeclareMathOperator{\Rea}{Re}
\newcommand{\areaAE}{\text{a.e.}}
\newcommand{\qe}{\text{q.e.}}
\title{A solution to the Erd\H{o}s Problem \#1040}
\author{Ioannis Tzachristas}
\date{Source: arXiv:2609.06050v1 (CC BY 4.0)}
\begin{document}
\maketitle

\noindent\emph{Source and licence.} A solution to the Erd\H{o}s Problem \#1040. Version arXiv:2609.06050v1 (CC BY 4.0), distributed by arXiv under the Creative Commons Attribution 4.0 International licence, http://creativecommons.org/licenses/by/4.0/, as recorded in the arXiv licence metadata for this exact version and shown on the abstract page https://arxiv.org/abs/2609.06050v1. Full licence notice: \texttt{licenses/arxiv-2609.06050-CC-BY-4.0.txt}.\par\smallskip

\noindent\textit{Editorial scope.} Counted unit: the article's lemma lem:realization (source0.tex lines 342--350). The article's complete original proof of it is delivered with it (source0.tex lines 352--422, one \texttt{proof} environment of 71 lines). No mathematics is written, repaired or re-derived: the statement and the proof are the article's own lines, in the article's own notation, and no retained line is substituted or re-flowed.  Retained uncounted as the source-local context the proof needs: lines 163--165; lines 210--213; lines 218--224.  Nothing was omitted from any retained span, so the package carries no omission record.  No retained line contains a citation, so the wrapper carries no bibliography and no bibliography entry.  No retained line contains a figure environment, an image-inclusion command or drawing code, so no image is delivered and the package declares no asset dependency.  Editorial formatting only, in addition: an AMS \texttt{amsart} preamble that restates the article's own declarations and macros used by the retained lines, namely the environments \texttt{lemma} on separate counters, together with the standard packages those lines need.  The retained environments print in this document as Lemma 1 in order of appearance, whereas the article prints the same items with the numbering of its own full document.  The complete native source of the article as distributed in the arXiv e-print for this exact version is bundled unchanged as \texttt{sources/209488/source0.tex}.
\begin{equation}\label{eq:log-int}
 \sup_{t\in K}\int_B |\log|z-t||\,dm(z)<\infty.
\end{equation}

\begin{equation}\label{eq:balayage-one}
 U_{\nu_w}(z)=\log|z-w|-U_\mu(w)
 \qquad\qe\text{ on }H.
\end{equation}

\begin{lemma}\label{lem:harnack}
Choose $0<r<R$ with $H\subset\{z:|z|<r\}$. For every $|w|=R$,
\begin{equation}\label{eq:domination}
 \nu_w\le C_{r,R}\mu,\qquad C_{r,R}=\frac{R+r}{R-r}.
\end{equation}
In particular, $R=2r$ gives $\nu_w\le3\mu$.
\end{lemma}

\begin{lemma}\label{lem:realization}
Let $h$ be a real harmonic polynomial with $\int h\,d\mu=0$.
There is a finite real signed measure $\sigma$ supported on $K$ such that
\begin{equation}\label{eq:realization}
 \sigma(K)=0,\qquad |\sigma|\le C_h\mu,\qquad
 U_\sigma=h\quad m\text{-}\areaAE\text{ on }H
\end{equation}
for some finite constant $C_h$.
\end{lemma}

\begin{proof}
Write $m_k=\int z^k\,d\mu(z)$. The centered condition allows us to
write
\begin{equation}\label{eq:h-expansion}
 h(z)=\Rea\sum_{k=1}^d a_k(z^k-m_k).
\end{equation}
If $h=0$, take $\sigma=0$. Otherwise choose $R$ larger than the
moduli of all points of $H$, and put $w_\theta=Re^{i\theta}$.
The uniformly convergent logarithmic expansion and
\eqref{eq:balayage-one} give, for each $\theta$,
\begin{equation}\label{eq:swept-expansion}
 U_{\nu_{w_\theta}}(z)
 =-\Rea\sum_{j=1}^{\infty}
       \frac{z^j-m_j}{jR^j}e^{-ij\theta}
 \quad m\text{-}\areaAE\text{ on }H.
\end{equation}
The series on the right converges uniformly for $(z,\theta)\in
H\times[0,2\pi]$.

Define the real trigonometric polynomial
\begin{equation}\label{eq:fourier-weight}
 b(\theta)=-2\sum_{k=1}^d kR^k\Rea(a_ke^{ik\theta}),
\end{equation}
and the real signed measure
\begin{equation}\label{eq:sigma}
 \sigma(A)=\int_0^{2\pi} b(\theta)\nu_{w_\theta}(A)
                          \frac{d\theta}{2\pi}.
\end{equation}
Harmonic measure is a Borel probability kernel in its pole, so this
integral defines a countably additive finite signed measure. For
example, countable additivity follows from dominated convergence
applied to disjoint unions, since $b$ is integrable and each
$\nu_{w_\theta}$ is a probability. Its total mass is zero because
$b$ has zero mean. Its support lies in $\partial H\subset K$.

By Lemma~\ref{lem:harnack}, for a constant $C_R$ independent of
$\theta$, one has $\nu_{w_\theta}\le C_R\mu$. Therefore
\begin{equation}\label{eq:variation-bound}
 |\sigma|(A)\le C_R\left(\int_0^{2\pi}|b(\theta)|
                       \frac{d\theta}{2\pi}\right)\mu(A).
\end{equation}
To see the total-variation inequality, first bound $|\sigma(A_j)|$
for each set in a finite Borel partition of $A$, sum, and then take
the supremum over such partitions.

We may integrate \eqref{eq:swept-expansion} against
$b(\theta)d\theta/(2\pi)$ and interchange the logarithmic integrations.
Indeed, for any bounded disk $B$ containing $H$,
\[
 \int_0^{2\pi}\!|b(\theta)|
 \int_B\!\int |\log|z-t||\,d\nu_{w_\theta}(t)\,dm(z)
 \frac{d\theta}{2\pi}<\infty
\]
by domination and \eqref{eq:log-int}. Fubini also deals with the
exceptional sets in \eqref{eq:swept-expansion}: they may depend on
$\theta$, but each has planar area zero. No common exceptional set
for all poles is required.

For $j,k\ge1$, Fourier orthogonality gives
\[
 \int_0^{2\pi}\Rea(a_ke^{ik\theta})
            \Rea(v_je^{-ij\theta})\frac{d\theta}{2\pi}
 =\begin{cases}
   \tfrac12\Rea(a_kv_k),&j=k,\\
   0,&j\ne k.
  \end{cases}
\]
Equations \eqref{eq:swept-expansion}--\eqref{eq:sigma} consequently
yield $U_\sigma=h$ almost everywhere on $H$. Together with
\eqref{eq:variation-bound}, this proves \eqref{eq:realization}.
\end{proof}

\end{document}
